% Exemplo: lista de exercícios com a classe exam (sem pontuação), com lacunas, verdadeiro ou falso
% e respostas que aparecem só na versão com a opção answers.
% Compilar: latexmk -pdf lista-exercicios.tex   Testado com: exam 2.7, TeX Live 2024 ou mais recente
\documentclass[12pt, a4paper]{exam}
\usepackage[T1]{fontenc}
\usepackage[brazilian]{babel}
\renewcommand{\solutiontitle}{\noindent\textbf{Resposta:}\enspace}
\qformat{\textbf{Exercício \thequestion}\hfill}   % título da questão em linha própria

\begin{document}

\begin{center}\large\bfseries Lista de exercícios 1: introdução ao \LaTeX\end{center}

%%BEGIN:listaexercicios
\begin{questions}
  \question O comando que inicia o corpo do documento é
    \fillin[\texttt{\textbackslash begin\{document\}}][5cm].

  \question Marque V (verdadeiro) ou F (falso).
  \begin{parts}
    \part \fillin[V][1cm] O \texttt{\%} inicia um comentário.
    \part \fillin[F][1cm] Uma linha em branco junta dois parágrafos.
  \end{parts}

  \question Escreva o código de uma lista numerada com dois itens.
  \begin{solutionorbox}[2.5cm]
    \texttt{\textbackslash begin\{enumerate\} \textbackslash item A
      \textbackslash item B
      \textbackslash end\{enumerate\}}
  \end{solutionorbox}
\end{questions}
%%END:listaexercicios

\end{document}
